% =====================================================================
% Lab Report 2 - Data Flow Diagram (DFD) Design
% Project: Nepali Lifestyle Based Diet and Fitness Advisory System
% Course: Software Engineering (Practical), ENCT 352
% =====================================================================
\documentclass[12pt,a4paper]{report}

% ---------------------------------------------------------------------
% PACKAGES
% ---------------------------------------------------------------------
\usepackage[a4paper,margin=1in]{geometry}   % Page size and 1-inch margins
\usepackage{newtxtext}                      % Times New Roman look-alike text font
\usepackage{newtxmath}                      % Matching math font
\usepackage{setspace}                       % Line spacing control
\usepackage{graphicx}                       % Image inclusion
\usepackage{float}                          % Precise figure placement with [H]
\usepackage{caption}                        % Caption formatting
\usepackage[hidelinks]{hyperref}            % Clickable TOC / cross-references
\usepackage{titlesec}                       % Section heading formatting
\usepackage{fancyhdr}                       % Header/footer and page numbering

% ---------------------------------------------------------------------
% GLOBAL FORMATTING
% ---------------------------------------------------------------------
\onehalfspacing                              % 1.5 line spacing throughout
\setlength{\parindent}{0pt}
\setlength{\parskip}{6pt}

% Justified paragraphs (default in LaTeX; \sloppy reduces overfull boxes)
\sloppy

% Caption style
\captionsetup{font=small,labelfont=bf,justification=centering}

% Section heading formatting (clean, professional numbering)
\titleformat{\chapter}[hang]{\normalfont\Large\bfseries}{\thechapter.}{10pt}{\Large}
\titlespacing*{\chapter}{0pt}{0pt}{20pt}
\titleformat{\section}{\normalfont\large\bfseries}{\thesection}{8pt}{}
\titleformat{\subsection}{\normalfont\bfseries}{\thesubsection}{6pt}{}

% Header/footer: page numbers only (no header rule), suppressed on title page
\pagestyle{fancy}
\fancyhf{}
\fancyfoot[C]{\thepage}
\renewcommand{\headrulewidth}{0pt}

% ---------------------------------------------------------------------
% DOCUMENT
% ---------------------------------------------------------------------
\begin{document}

% =====================================================================
% TITLE PAGE (no page number)
% =====================================================================
\begin{titlepage}
    \centering
    \thispagestyle{empty}

    \vspace*{0.3cm}
    \includegraphics[width=0.85\textwidth]{tu_banner.png}\\[2cm]

    {\Large \textbf{Lab Report 2}}\\[0.5cm]
    {\large \textbf{Data Flow Diagram (DFD) Design}}\\[1.5cm]

    {\normalsize Project:}\\[3pt]
    {\large \textbf{Nepali Lifestyle Based Diet and Fitness Advisory System}}\\[2cm]

    \begin{tabular}{ll}
        \textbf{Course} & Software Engineering (Practical) \\[4pt]
        \textbf{Course Code} & ENCT 352 \\
    \end{tabular}

    \vspace{2cm}

    {\normalsize \textbf{Submitted By}}\\[8pt]
    \begin{tabular}{ll}
        HIMAL BHANDARI   & PUR080BCT034 \\
        KSHITIZ PORTEL    & PUR080BCT039 \\
        MANAV ACHARYA     & PUR080BCT041 \\
        MAUSAM PARAJULI   & PUR080BCT045 \\
    \end{tabular}

    \vfill

    \begin{tabular}{ll}
        \textbf{Submitted To} & Department of Electronics and Computer Engineering \\
        \textbf{Instructor} & Asst. Prof. Binay Lal Shrestha \\
        \textbf{Date} & July 06, 2026 \\
    \end{tabular}

\end{titlepage}

% =====================================================================
% TABLE OF CONTENTS
% =====================================================================
\pagenumbering{roman}
\tableofcontents
\newpage
\pagenumbering{arabic}
\setcounter{page}{1}

% =====================================================================
% CHAPTER: INTRODUCTION
% =====================================================================
\chapter{Introduction}

A Data Flow Diagram (DFD) is a graphical representation that shows how data moves through a
system, identifying the external entities that interact with the system, the processes that
transform data, the data stores where information is retained, and the flows that connect them.
Unlike the Software Requirements Specification, which describes \emph{what} a system must do
in textual form, a DFD models \emph{how} data actually travels between the user and the
internal processes, giving developers a clearer picture of system behaviour before design and
coding begin.

DFDs are prepared after the Software Requirements Specification because the requirements first
establish the functional scope of the system, and the DFD then translates that scope into a
structured, visual data flow model. This model helps the development team identify the exact
inputs, outputs, and internal processing steps required to satisfy the requirements, and it
exposes any gaps or inconsistencies before implementation starts.

For the Nepali Lifestyle Based Diet and Fitness Advisory System, the DFDs prepared in this lab
illustrate how user data, food images, and health goals flow between the user, the internal
processes such as user management, food analysis, and recommendation generation, and the
supporting data stores. These diagrams form the structural basis for the system design that
follows in later phases of the project.

% =====================================================================
% CHAPTER: DATA FLOW DIAGRAMS
% =====================================================================
\chapter{Data Flow Diagrams}

This chapter presents the Context Diagram, the Level 1 DFD, and the Level 2 DFD developed for
the Nepali Lifestyle Based Diet and Fitness Advisory System. Each diagram is described in terms
of its external entities, processes, data stores, and data flows.

% ---------------------------------------------------------------------
\section{Context Diagram}

\begin{figure}[H]
    \centering
    \includegraphics[width=0.85\textwidth]{clear_context_diagram.png}
    \caption{Context Diagram of the Nepali Lifestyle Based Diet and Fitness Advisory System}
    \label{fig:context}
\end{figure}

The Context Diagram shown in Figure~\ref{fig:context} represents the highest level of the DFD,
labelled as process 0. It provides an overview of the entire system as a single process and
defines its boundary with the external environment.

The only external entity in this diagram is the \textbf{User}. Two data flows connect the User
to the Diet and Fitness Advisory System: ``Profile setup and meal photos'', which
carries the user's profile information and meal images into the system, and
``Diet and Fitness feedback'', which carries the generated recommendations back to the
User. No data stores are shown at this level, since the Context Diagram treats the system as a
single, undivided process.

As the first and highest-level diagram, the Context Diagram does not expand any previous level.
It instead establishes the overall system boundary that the Level 1 DFD subsequently decomposes
into individual processes.

% ---------------------------------------------------------------------
\section{Level 1 DFD}

\begin{figure}[H]
    \centering
    \includegraphics[width=0.85\textwidth]{clear_level1_dfd.png}
    \caption{Level 1 DFD of the Nepali Lifestyle Based Diet and Fitness Advisory System}
    \label{fig:level1}
\end{figure}

Figure~\ref{fig:level1} presents the Level 1 DFD, which decomposes the single process shown in
the Context Diagram into three major processes of the system.

The external entity remains the \textbf{User}. The three processes at this level are
\textbf{1 User Management}, \textbf{2 Food Analysis}, and \textbf{3 Recommendation Generation}.
Three data stores support these processes: \textbf{D1 User DB}, \textbf{D2 Food DB}, and
\textbf{D3 Report DB}. The User provides personal information to User Management, which stores
and retrieves user data from the User DB. The User also provides food details or a meal photo
to Food Analysis, which stores food information in the Food DB. Recommendation Generation
retrieves food information from the Food DB, stores and retrieves reports from the Report DB,
and finally sends the diet and fitness plan along with the nutrition report back to the User.

This diagram expands the single process of the Context Diagram into the three functional areas
of the system: managing user profiles, analysing food, and generating recommendations. Each of
these processes now has a defined data store and its own set of data flows, giving a clearer
view of how the system separates its responsibilities internally.

% ---------------------------------------------------------------------
\section{Level 2 DFD}

\begin{figure}[H]
    \centering
    \includegraphics[width=0.85\textwidth]{clear_level2_dfd.png}
    \caption{Level 2 DFD - Decomposition of Process 3 (Recommendation Generation)}
    \label{fig:level2}
\end{figure}

Figure~\ref{fig:level2} shows the Level 2 DFD, which further decomposes process 3,
Recommendation Generation, from the Level 1 DFD into four sub-processes.

The external entity is again the \textbf{User}. The four sub-processes are
\textbf{3.1 Read User Data}, \textbf{3.2 Analyze Nutrition}, \textbf{3.3 Generate Diet and
Fitness Plan}, and \textbf{3.4 Prepare Report}. Two data stores appear at this level, the
\textbf{Food DB} and the \textbf{Report DB}. The User provides health goals to Read User Data,
which passes this information to Analyze Nutrition. Analyze Nutrition stores food information in
the Food DB, which is later retrieved by Generate Diet and Fitness Plan. The output of this
sub-process flows into Prepare Report, which stores and retrieves report data from the Report
DB and finally provides the diet plan, fitness plan, and health report back to the User.

This diagram refines process 3 of the Level 1 DFD by breaking down Recommendation Generation
into a sequential set of sub-processes that read user health goals, analyse nutrition against
stored food information, generate the actual plan, and prepare the final report. It shows the
internal working of the recommendation process without altering the data flows already defined
at the higher level.

% =====================================================================
% CHAPTER: OVERALL DISCUSSION
% =====================================================================
\chapter{Overall Discussion}

The three diagrams developed in this lab show a clear progression from a general overview to a
detailed internal structure. The Context Diagram treats the Nepali Lifestyle Based Diet and
Fitness Advisory System as a single process interacting with the User, establishing the system
boundary. The Level 1 DFD then breaks this single process into three functional areas, User
Management, Food Analysis, and Recommendation Generation, each supported by its own data store.
The Level 2 DFD goes further by decomposing Recommendation Generation into four sequential
sub-processes, revealing how user health goals are read, analysed against stored food data, and
converted into a diet and fitness plan.

Across these levels, information consistently flows from the User into the system, through the
relevant process, into or out of a data store, and back to the User as feedback. This layered
structure mirrors the system architecture described in the SRS, where profile management, food
analysis, and recommendation generation form distinct but connected modules built on a
relational data layer.

These diagrams directly support the functional requirements defined in Lab 1, such as profile
management, meal logging, and personalized advisory generation. By visualizing the exact data
flows before implementation, the development team can verify that every requirement has a
corresponding process and data path, making the DFDs a valuable reference during the design and
coding phases that follow.

% =====================================================================
% CHAPTER: CONCLUSION
% =====================================================================
\chapter{Conclusion}

In this lab, the Context Diagram, Level 1 DFD, and Level 2 DFD were developed for the Nepali
Lifestyle Based Diet and Fitness Advisory System, based directly on the requirements identified
in Lab 1. These diagrams illustrate how user information, meal data, and health goals move
through the system's processes and data stores at increasing levels of detail.

Data Flow Diagrams are important because they translate textual requirements into a visual model
of system behaviour, making it easier to identify processes, data stores, and flows that must be
implemented. Working across multiple levels, from the Context Diagram down to the Level 2
decomposition, allows the team to understand both the overall system boundary and the detailed
internal workings of specific processes such as recommendation generation.

Completing this lab has clarified how the major modules of the system interact and exchange
data, providing a solid foundation for the system design phase. Lab 2 therefore prepares the
project for the next phase of development, in which these data flows will be translated into
concrete database schemas and module-level designs.

\end{document}
